% ===== 问题三正文 =====
% 本文件汇总第三问；各部分正文、公式与图表请编辑 sections/q3/ 下的对应子文件。
\subsection{问题三模型建立与求解}\label{sec:q3}

第二问给出了已发现单源的主动选点方法。面对多个数量未知的源，还需要安排覆盖扫描，并决定先处理哪个任务。本文将待扫描点与已发现源合并为动态任务集合，以当前定位区域表示源的位置，每次选择一项任务执行，再根据新观测重新安排路线。

% 第三问辅助示意图。颜色沿用 q1_case.tex，局部由 sections/q3.tex 引入。
% 图宽按 16 cm 设计；几何图按注释给定比例，概念规划图不代表实际轨迹。
% 仅使用已有的 arrows.meta、calc、positioning、shapes.geometric、shapes.misc 库。

% ---------- 原点扫描、七环布局及最近站的覆盖保证 ----------
\newcommand{\QThreeCoverageFigure}{\resizebox{\linewidth}{!}{%
\begin{tikzpicture}[>=Stealth,line cap=round,line join=round,
  font=\fontsize{10}{12}\selectfont,text=qoneink,
  dot/.style={circle,inner sep=1.7pt,fill=qoneink},
  site/.style={rectangle,inner sep=2.2pt,fill=qonecover,draw=qonecover!80!black},
  info/.style={fill=none,inner sep=1.5pt,align=center}]
\path[use as bounding box] (0,0) rectangle (16,8.05);
\fill[white] (0,0) rectangle (16,8.05);
\draw[black!15] (7.7,.8)--(7.7,7.65);
\node[font=\bfseries] at (3.7,7.65) {(a) 原点扫描与七个环上检测点};
\node[font=\bfseries] at (11.9,7.65) {(b) 每个可能位置附近都有检测点};

% 左图比例：1 cm = 600 m，三类圆都用同一比例。
\coordinate (O) at (3.7,4.15);
\fill[qoneblue!5] (O) circle[radius=3];
\begin{scope}
  \clip (O) circle[radius=3];
  \foreach \k in {0,...,6} {
    \pgfmathsetmacro{\ang}{360*\k/7}
    \draw[qonecover!60,line width=.45pt] ($(O)+(\ang:1.665)$) circle[radius=1.666667];
  }
  \fill[qonecover!16] ($(O)+(1.665,0)$) circle[radius=1.666667];
  \draw[qonecover,line width=.9pt] ($(O)+(1.665,0)$) circle[radius=1.666667];
\end{scope}
\draw[qoneblue,line width=.9pt] (O) circle[radius=3];
\draw[qoneink!55,densely dashed] (O) circle[radius=1.665];
\foreach \k in {0,...,6} {
  \pgfmathsetmacro{\ang}{360*\k/7}
  \node[site] at ($(O)+(\ang:1.665)$) {};
}
\node[dot] at (O) {};
\node[below right,info] at ($(O)+(.18,-.18)$) {原点 $O$};
\draw[<->,qoneink!75] (O)--($(O)+(235:1.665)$);
\node[info,anchor=west] (ringlabel) at (.2,1.65) {$999$ m};
\draw[qoneink!55,line width=.5pt] (ringlabel.north east)--(2.2,2.25)--($(O)+(235:1.665)$);
\draw[<->,qonecover] ($(O)+(1.665,0)$)--($(O)+(1.665,0)+(65:1.666667)$);
\node[info,text=qonecover] (radiuslabel) at (6.1,6.8) {$1000$ m};
\draw[qonecover,line width=.5pt] (radiuslabel.south)--($(O)+(1.665,0)+(65:1.666667)$);
\node[info,text=qoneblue] at (3.7,.8) {场地半径 $1800$ m};
\node[info,font=\fontsize{9.5}{11.5}\selectfont] at (3.7,.15)
  {七个检测点覆盖场地，原点先获取信息};

% 右图仍用 1 cm = 600 m；画出角度最不利的边界点。
\coordinate (C) at (9.55,4.0);
\coordinate (s) at ($(C)+(1.665,0)$);
\coordinate (p) at ($(C)+(25.714286:3)$);
\fill[qoneblue!7] (C)--($(C)+(-25.714286:3)$)
  arc[start angle=-25.714286,end angle=25.714286,radius=3]--cycle;
\draw[qoneblue,line width=.9pt] ($(C)+(-25.714286:3)$)
  arc[start angle=-25.714286,end angle=25.714286,radius=3];
\draw[qoneblue!55,densely dotted] (C)--($(C)+(-25.714286:3)$);
\draw[qoneink!65,densely dashed] (C)--($(C)+(3.25,0)$);
\draw[qoneink,line width=.7pt] (C)--(p);
\draw[qonecover,line width=1pt] (s)--(p);
\draw[qoneorange,line width=.8pt] ($(C)+(.75,0)$)
  arc[start angle=0,end angle=25.714286,radius=.75];
\node[text=qoneorange,info] at ($(C)+(.55,.82)$) {$\alpha\le\pi/7$};
\node[dot] at (C) {};
\node[site] at (s) {};
\node[dot,fill=qonefail] at (p) {};
\node[below left,info] at (C) {$O$};
\node[below,info] at ($(s)+(0,-.1)$) {检测点 $x_k$};
\node[above right,info] at ($(p)+(.10,.18)$) {源的可能位置 $p$};
\node[info] at (10.4,5.7) {$r\le1800$ m};
\node[anchor=west,info,text=qonecover] (distlabel) at (13.05,4.9) {$d\le1000$ m};
\draw[qonecover,line width=.5pt] (distlabel.west)--($(s)!.6!(p)+(.08,0)$);
\node[info,align=left,anchor=north west] at (8.25,2.50)
  {$d^2=r^2+999^2-2\times999r\cos\alpha$\\[3pt]
   关于 $r$ 为凸函数，只需检查 $r=0,1800$\\[3pt]
   最不利情形仍有 $d\le999<1000$ m};
\node[info,font=\fontsize{9.5}{11.5}\selectfont] at (11.8,.15)
  {按实际比例，场地外的接收圆弧省略};
\end{tikzpicture}
}}

% ---------- 将覆盖任务和已知源任务合并，并在真实反馈后重新规划 ----------
\newcommand{\QThreePlanningFigure}{\resizebox{\linewidth}{!}{%
\begin{tikzpicture}[>=Stealth,line cap=round,line join=round,
  font=\fontsize{10}{12}\selectfont,text=qoneink,
  dot/.style={circle,inner sep=1.8pt,fill=qoneink},
  past/.style={circle,inner sep=2.1pt,draw=qoneink!65,fill=white,line width=.8pt},
  site/.style={rectangle,inner sep=2.2pt,fill=qonecover,draw=qonecover!80!black},
  source/.style={circle,inner sep=1.7pt,fill=qoneblue},
  info/.style={fill=white,inner sep=1.5pt,align=center}]
\path[use as bounding box] (0,0) rectangle (16,7.45);
\fill[white] (0,0) rectangle (16,7.45);
\draw[black!15] (8,1.5)--(8,7.05);
\node[font=\bfseries] at (4,7.0) {(a) 合并任务，安排访问顺序};
\node[font=\bfseries] at (12,7.0) {(b) 到达后读取检测结果，再安排任务};

\begin{scope}[shift={(.35,1.5)}]
  \coordinate (x) at (.8,2.0);
  \coordinate (old) at (1.0,.55);
  \coordinate (F1) at (2.0,4.45);
  \coordinate (F2) at (6.0,1.0);
  \coordinate (z1) at (5.15,4.05);
  \coordinate (z2) at (3.35,.55);
  \foreach \pt in {z1,z2} {
    \filldraw[fill=qoneblue!10,draw=qoneblue!60,line width=.6pt]
      ($(\pt)+(-.65,-.1)$)--($(\pt)+(-.3,-.3)$)--($(\pt)+(.58,-.17)$)
      --($(\pt)+(.7,.08)$)--($(\pt)+(.25,.27)$)--($(\pt)+(-.5,.2)$)--cycle;
    \node[source] at (\pt) {};
  }
  \draw[->,qoneink!65,densely dashed,line width=.7pt] (x)--(F1);
  \draw[->,qoneink!65,densely dashed,line width=.7pt] (F1)--(z1);
  \draw[->,qoneink!65,densely dashed,line width=.7pt] (z1)--(F2);
  \draw[->,qoneink!65,densely dashed,line width=.7pt] (F2)--(z2);
  \draw[->,qoneorange,line width=1.5pt] (x)--(F1)
    node[midway,left,info,text=qoneorange] {这一步\\只去这里};
  \node[dot] at (x) {};
  \node[past] at (old) {};
  \node[site] at (F1) {};
  \node[site] at (F2) {};
  \node[below left,info] at (x) {$x_t$};
  \node[below,info] at (old) {$P_1$};
  \node[above,info] at (F1) {$F_1$：下一个扫描点};
  \node[right,info] at (F2) {$F_2$};
  \node[above,info] at ($(z1)+(0,.30)$) {定位区域圆心 $z_1$};
  \node[below,info] at ($(z2)+(0,-.30)$) {定位区域圆心 $z_2$};
\end{scope}

\begin{scope}[shift={(8.35,1.5)}]
  \coordinate (old) at (1.0,.55);
  \coordinate (F1) at (2.0,4.45);
  \coordinate (F2) at (6.0,1.0);
  \coordinate (z1) at (5.15,4.05);
  \coordinate (z2) at (3.35,.55);
  \coordinate (z3) at (2.7,2.4);
  \foreach \pt in {z1,z2,z3} {
    \filldraw[fill=qoneblue!10,draw=qoneblue!60,line width=.6pt]
      ($(\pt)+(-.65,-.1)$)--($(\pt)+(-.3,-.3)$)--($(\pt)+(.58,-.17)$)
      --($(\pt)+(.7,.08)$)--($(\pt)+(.25,.27)$)--($(\pt)+(-.5,.2)$)--cycle;
    \node[source] at (\pt) {};
  }
  \draw[->,qoneink!65,densely dashed,line width=.7pt] (F1)--(z3);
  \draw[->,qoneink!65,densely dashed,line width=.7pt] (z3)--(z2);
  \draw[->,qoneink!65,densely dashed,line width=.7pt] (z2)--(F2);
  \draw[->,qoneink!65,densely dashed,line width=.7pt] (F2)--(z1);
  \node[past] at (old) {};
  \node[past] at (F1) {};
  \draw[qoneorange,line width=1.1pt] (F1) circle[radius=.17];
  \node[site] at (F2) {};
  \node[below,info] at (old) {$P_1$};
  \node[above,info] at (F1) {$F_1$ 扫描完成，加入 $P_t$};
  \node[right,info] at (F2) {$F_2$};
  \node[above,info] at ($(z1)+(0,.30)$) {$z_1$};
  \node[below,info] at ($(z2)+(0,-.30)$) {$z_2$};
  \node[anchor=west,info,text=qoneblue] at ($(z3)+(.6,.3)$) {新源的区域圆心 $z_3$};
  \node[anchor=east,info,text=qoneorange] at ($(F1)+(-.3,-.38)$) {当前位置};
\end{scope}

\node[past] at (.65,.88) {};
\node[anchor=west] at (.88,.88) {已扫描位置 $P_t$};
\node[site] at (4.1,.88) {};
\node[anchor=west] at (4.35,.88) {待扫描的点 $F_t$};
\node[source] at (8.1,.88) {};
\node[anchor=west] at (8.35,.88) {定位区域圆心};
\draw[qoneink!65,densely dashed] (12.0,.88)--(12.65,.88);
\node[anchor=west] at (12.85,.88) {当前计划路线};
\end{tikzpicture}
}}

% ---------- 按包围圆尺度生成候选测点，再作接收距离筛选 ----------
\newcommand{\QThreeProbeFigure}{\resizebox{\linewidth}{!}{%
\begin{tikzpicture}[>=Stealth,line cap=round,line join=round,
  font=\fontsize{10.5}{13}\selectfont,text=qoneink,
  cand/.style={circle,inner sep=2.4pt,fill=qoneblue,draw=white,line width=.7pt},
  reject/.style={cross out,draw=qonefail,line width=1.0pt,inner sep=2.6pt},
  lab/.style={fill=none,inner sep=2pt,align=left},
  lead/.style={line width=.65pt,draw=qoneink!60}]
\path[use as bounding box] (0,-.35) rectangle (16,9.2);
\fill[white] (0,-.35) rectangle (16,9.2);

% 教学几何，等比例绘制：1 cm=200 m。
% x=(0,0), C=[200,1400]×[-20,20], z=(800,0), r=sqrt(600^2+20^2)。
% 距离域边界：(|q_x-800|+600)^2+(|q_y|+20)^2=999.9^2。
% 填充由四个圆盘的真实交集给出，不用椭圆替代。
\pgfmathsetmacro{\rr}{sqrt(3*3+.1*.1)}
\pgfmathsetmacro{\receive}{4.9995}
\coordinate (z) at (8,4.7);
\coordinate (x) at (4,4.7);
\coordinate (v1) at (5,4.6);
\coordinate (v2) at (11,4.6);
\coordinate (v3) at (11,4.8);
\coordinate (v4) at (5,4.8);

% 浅蓝域：候选点到每个顶点的距离均小于999.9 m。
\begin{scope}
  \foreach \v in {v1,v2,v3,v4} \clip (\v) circle[radius=\receive];
  \fill[qoneblue!9] (0,0) rectangle (16,9.2);
  \foreach \v in {v1,v2,v3,v4}
    \draw[qoneblue!65,densely dashed,line width=.7pt]
      (\v) circle[radius=\receive];
\end{scope}

% 包围圆只用于确定布点尺度，灰紫虚线与接收域边界区分。
\draw[qonepurple!75,dash pattern=on 4pt off 3pt,line width=.8pt]
  (z) circle[radius=\rr];

% 淡网格强调横纵各五个采样位置。
\foreach \a in {-.7,-.35,0,.35,.7}
  \draw[black!12,line width=.35pt]
    ($(z)+(\a*\rr,-.4*\rr)$)--($(z)+(\a*\rr,.4*\rr)$);
\foreach \b in {-.4,-.15,0,.15,.4}
  \draw[black!12,line width=.35pt]
    ($(z)+(-.7*\rr,\b*\rr)$)--($(z)+(.7*\rr,\b*\rr)$);

\filldraw[fill=qoneregion!25,draw=qoneregion,line width=.9pt]
  (v1)--(v2)--(v3)--(v4)--cycle;
\foreach \v in {v1,v2,v3,v4} \fill[qoneregion] (\v) circle[radius=.035];

% 半径标记避开下排点，不与图外标注交叉。
\draw[->,qonepurple!80,line width=.7pt]
  (z)--($(z)+(-62:\rr)$);
\node[lab,text=qonepurple] at ($(z)+(.78,-1.88)$) {$r$};

\foreach \a in {-.7,-.35,0,.35,.7} {
  \foreach \b in {-.4,-.15,0,.15,.4} {
    \pgfmathtruncatemacro{\bad}{abs(\a)>.6}
    \ifnum\bad=1
      \node[reject] at ($(z)+(\a*\rr,\b*\rr)$) {};
    \else
      \node[cand] at ($(z)+(\a*\rr,\b*\rr)$) {};
    \fi
  }
}
\node[reject] at (x) {};
\node[lab,anchor=north] at ($(x)+(0,-.24)$) {当前位置 $x$};
\node[lab,anchor=north west,inner sep=1pt] at ($(z)+(.22,.94)$) {$z$};

% 直接标注布置在数据区外。
\node[lab,anchor=east,text=qonepurple] (circlelabel) at (4.75,7.55) {包围圆};
\draw[lead,qonepurple!70] (circlelabel.east)--(5.1,7.55)--($(z)+(133:\rr)$);
\node[lab,anchor=west,text=qoneblue] (domainlabel) at (11.4,8.0)
  {可靠接收区};
\draw[lead,qoneblue!70] (domainlabel.west)--(10.65,8.0)--(8.71,8.0);
\node[lab,anchor=west,text=qoneregion] (regionlabel) at (12.0,5.65)
  {当前定位区域 $\widehat C_f$};
\draw[lead,qoneregion] (regionlabel.west)--(11.65,5.65)--(10.82,4.8);

% 局部方向轴只解释布点方向，不作为数值坐标轴。
\draw[->,qoneink!75,line width=.8pt] (2.2,1.45)--(4.35,1.45)
  node[right,inner sep=2pt] {$e$};
\draw[->,qoneink!75,line width=.8pt] (2.2,1.45)--(2.2,2.75)
  node[above,inner sep=2pt] {$e_\perp$};
\node[anchor=north,align=center,font=\fontsize{9.5}{12}\selectfont]
  at (3.3,1.2) {$e$ 指向包围圆心};

% 图例使用颜色与形状双重编码。
\node[cand] at (3.2,.15) {};
\node[anchor=west] at (3.45,.15) {通过距离检验：15点};
\node[reject] at (9.1,.15) {};
\node[anchor=west] at (9.35,.15) {未通过：10点及当前位置};
\end{tikzpicture}
}}

% ---------- 与算法3一致的滚动求解流程 ----------
\newcommand{\QThreeFlowFigure}{\resizebox{\linewidth}{!}{%
\begin{tikzpicture}[paper flow]
\path[use as bounding box] (0,.1) rectangle (16,11.9);
\PaperFlowCard{start}{(8,11.0)}{6.6}{1.4}{Blue}
  {初始化与扫描}{原点扫描20个频道，布置七个检测点}

% 完成条件与图1一致；短分支结束，未完成时向下执行。
\node[diamond,aspect=4.5,draw=flowBlueTitle,fill=flowBlueBody,
  line width=.8pt,inner sep=0pt,text width=35mm,align=center,
  minimum width=60mm,minimum height=14mm] (done) at (8,9.1)
  {未知频道均已排除\\已发现源全部清除？};
\PaperFlowCard{finish}{(2.15,9.1)}{3.2}{1.4}{Green}
  {任务结束}{核对清除记录}
\node[flow label,text=flowNote,text width=42mm,
  font=\linespread{1}\fontsize{9}{11.5}\selectfont] at (2.15,7.72)
  {未知频道排除依据\\全场覆盖或已发现16个源};

\PaperFlowCard{plan}{(8,6.9)}{6.6}{1.65}{Warm}
  {构造任务并选择下一步}
  {先清除当前位置可清除源\\调整扫描点，合并任务并用2-opt排序}

% 任务类型由左右分支表示，不再另画一个重复说明的判断框。
\PaperFlowCard{coverage}{(4,4.6)}{6.0}{1.65}{Blue}
  {覆盖扫描}{到站扫描未知频道\\按需补测已发现源}
\PaperFlowCard{source}{(12,4.6)}{6.0}{1.65}{Blue}
  {主动测向或清除}
  {满足条件则清除，否则选点检测\\必要时在包围圆心测向}
\PaperFlowCard{update}{(8,2.5)}{6.6}{1.65}{Blue}
  {真实反馈与状态更新}
  {按需追加扫描或测向\\更新定位区域、覆盖记录与清除状态}

\draw[flow arrow] (start)--(done);
\draw[flow arrow] (done.west)--node[flow label,above]{是}(finish.east);
\draw[flow arrow] (done.south)--node[flow label,right]{否}(plan.north);
\draw[flowBlueTitle,line width=.8pt] (plan.south)--(8,5.72);
\draw[flow arrow] (8,5.72)--(4,5.72)--(coverage.north);
\draw[flow arrow] (8,5.72)--(12,5.72)--(source.north);
\node[flow label,above] at (3.7,5.72) {扫描任务};
\node[flow label,above] at (12.3,5.72) {源任务};
\draw[flow arrow] (coverage.south)--(4,2.5)--(update.west);
\draw[flow arrow] (source.south)--(12,2.5)--(update.east);
\draw[flow arrow] (update.south)--(8,1.0)--(15.65,1.0)
  --(15.65,9.1)--(done.east);
\node[flow label,text=flowBlueTitle] at (11.5,.53)
  {每执行一步，按观测结果重新规划};
\end{tikzpicture}%
}}
%
% 第三问圆心测向：构造扇形覆盖圆，并绘制保守半径上界。
% 左图半角放大至18度，整体旋转20度；右图计算仍用1.005度。
% 示意角度与理论数据的说明放在正文图注中。
\newcommand{\QThreeContractionFigure}{\resizebox{\linewidth}{!}{%
\begin{tikzpicture}[>=Stealth,line cap=round,line join=round,
  font=\fontsize{10}{12.5}\selectfont,text=qoneink,
  lab/.style={fill=none,inner sep=2pt,align=center},
  lead/.style={line width=.65pt},
  point/.style={circle,inner sep=1.7pt,fill=qoneink}]
\path[use as bounding box] (0,0) rectangle (16.6,8.0);
\fill[white] (0,0) rectangle (16.6,8.0);
\draw[black!12] (7.9,1.0)--(7.9,7.30);
\node[font=\bfseries] at (3.85,7.6) {(a) 在圆心再次测量，缩小定位区域};
\node[font=\bfseries] at (12.15,7.6) {(b) 多次测量后的半径上界};

% ---------- 左：原包围圆、误差扇形及扇形覆盖圆 ----------
\pgfmathsetmacro{\rad}{2.3}
\pgfmathsetmacro{\coverR}{\rad/(2*cos(18))}
\coordinate (z) at (3.00,4.15);
\coordinate (c) at ($(z)+(20:\coverR)$);
\coordinate (a) at ($(z)+(2:\rad)$);
\coordinate (b) at ($(z)+(38:\rad)$);

\fill[qoneblue!5] (z) circle[radius=\rad];
\draw[qoneblue!85,dash pattern=on 3.5pt off 2.5pt,line width=.85pt]
  (z) circle[radius=\rad];
\fill[qonepurple!7] (c) circle[radius=\coverR];
\fill[qoneorange!27] (z)--(a)
  arc[start angle=2,end angle=38,radius=\rad]--cycle;
\draw[qonepurple,line width=1.0pt] (c) circle[radius=\coverR];
\draw[qoneorange,line width=1.05pt] (a)--(z)--(b);
\draw[qoneorange,line width=1.05pt] (a)
  arc[start angle=2,end angle=38,radius=\rad];

\draw[qoneorange,line width=.65pt] ($(z)+(2:.55)$)
  arc[start angle=2,end angle=38,radius=.55];
\node[text=qoneorange,inner sep=1pt] at ($(z)+(.89,.31)$) {$2\bar\varepsilon$};
\draw[->,qoneblue!85,line width=.7pt] (z)--($(z)+(140:\rad)$);
\node[lab,text=qoneblue] at ($(z)+(-1.12,.54)$) {$r$};
\node[point] at (z) {};
\node[lab,anchor=north east] at ($(z)+(-.10,-.16)$) {检测点 $z$};
\foreach \p in {a,b} \fill[qoneorange] (\p) circle[radius=.042];

\node[lab,text=qoneblue] at (2.10,6.85) {原包围圆 $B(z,r)$};
\node[lab,text=qoneorange,anchor=west,align=left] (sectorlabel) at (5.70,6.20)
  {新区域位于\\此扇形内};
\draw[lead,qoneorange!85] (sectorlabel.south west)--(5.70,5.68)--(5.18,4.93);
\node[lab,text=qonepurple,anchor=west] (coverlabel) at (5.28,2.55) {能罩住扇形的圆};
\draw[lead,qonepurple] (coverlabel.west)--(5.02,2.55)--($(c)+(-49:\coverR)$);
\node[text=qonepurple,font=\fontsize{11}{14}\selectfont] at (3.85,1.13)
  {$\displaystyle r_{\rm new}\le R=\frac{r}{2\cos\bar\varepsilon}$};

% ---------- 右：线性坐标，递推数据未作平滑或修改 ----------
\coordinate (origin) at (8.95,1.55);
\foreach \v in {40,80,120,160} {
  \draw[black!12,line width=.4pt] (8.95,{1.55+\v*.029})--(16.05,{1.55+\v*.029});
}
\draw[->,qoneink!65,line width=.7pt] (origin)--(16.20,1.55);
\draw[->,qoneink!65,line width=.7pt] (origin)--(8.95,6.85);
\node[anchor=west] at (9.02,6.94) {半径上界 / m};
\node at (12.55,.42) {圆心处的检测次数 $k$};
\foreach \v in {0,40,80,120,160} {
  \node[anchor=east,font=\fontsize{9.5}{12}\selectfont]
    at (8.80,{1.55+\v*.029}) {\v};
}

\pgfmathsetmacro{\thresholdY}{1.55+19.5*.029}
\draw[qonefail,dash pattern=on 3pt off 2pt,line width=.8pt]
  (8.95,\thresholdY)--(16.05,\thresholdY);
\node[lab,anchor=south west,text=qonefail,font=\fontsize{9}{11}\selectfont]
  at (9.15,{\thresholdY+.07}) {清除所需半径：$19.5$ m};

\foreach \k in {0,...,4} {
  \pgfmathsetmacro{\bound}{160/(2*cos(1.005))^\k}
  \coordinate (p\k) at ({9.15+\k*1.57},{1.55+\bound*.029});
  \draw[qoneink!50,line width=.5pt] ({9.15+\k*1.57},1.55)--({9.15+\k*1.57},1.44);
  \node[anchor=north,font=\fontsize{9.5}{12}\selectfont]
    at ({9.15+\k*1.57},1.31) {\k};
}
\draw[qoneblue,line width=1.0pt] (p0)--(p1)--(p2)--(p3)--(p4);
\foreach \k in {0,1,2,3}
  \node[point,fill=qoneblue,draw=white,line width=.5pt,inner sep=2.2pt] at (p\k) {};
\node[point,fill=qonecover,draw=white,line width=.5pt,inner sep=2.5pt] at (p4) {};

\node[lab,anchor=south west,text=qoneblue] at ($(p0)+(.02,.05)$) {$160.00$};
\node[lab,anchor=south west,text=qoneblue] at ($(p1)+(.03,.06)$) {$80.01$};
\node[lab,anchor=south west,text=qoneblue] at ($(p2)+(.03,.08)$) {$40.01$};
\node[lab,anchor=south,text=qoneblue] at (13.86,3.10) {$20.01$};
% 末端两点靠近门限，数值放在曲线上方，用短引线对应，不占横轴空间。
\draw[lead,qoneblue!70] ($(p3)+(0,.04)$)--(13.86,3.09);
\node[lab,anchor=south,text=qonecover] at (15.46,3.10) {$10.01$};
\draw[lead,qonecover!80] ($(p4)+(0,.05)$)--(15.43,3.09);
\end{tikzpicture}%
}}
%
\Needspace{30mm}
\subsubsection{模型准备}\label{sec:q3-1}

\textbf{（1）已知条件与求解难点。}第三问要求在半径1800 m的圆形场地 $\Omega$ 内发现并清除全部全向静止干扰源。设备共有20个频道，干扰源最多有16个；每个源的接收半径 $R_f$ 在一次任务中固定，且满足 $1000\le R_f\le1500$ m。示向度的测量误差为 $\pm1^\circ$。机器狗与源的距离不超过5 m时，设备会给出近场反馈，即提示源就在5 m以内；距离不超过20 m时，可以进行光学精确定位和清除。源的位置、实际数量和接收半径均不能预先读取。

与第二问相比，本问的困难有三点：一是尚未发现的源没有定位区域，需要逐个检测未知频道来寻找；二是新的检测结果会改变已发现源的定位区域，也可能增加新的任务，因此出发前无法确定完整路线；三是清除了全部已发现源之后，还需判断是否存在遗漏。因此，搜索、定位和清除要根据每一步的实际检测结果一起安排。

\textbf{（2）符号约定。}记机器狗位置为 $x_t$，当前检测频道为 $f_t$，按题目规则累计的任务耗时为 $T_t$；已发现源和已清除源所在的频道集合分别为 $\mathcal D_t$、$\mathcal K_t$，尚不能确定是否有源的频道集合为 $\mathcal U_t$，已确认没有源的频道集合为 $\mathcal A_t$。将当前已知信息统一记为
\begin{equation}\label{eq:q3-1}
S_t=\left(x_t,f_t,T_t,\mathcal U_t,\mathcal A_t,\mathcal D_t,\mathcal K_t,
\{\widehat C_f,\mathcal O_f^+,\mathcal O_f^-\},P_t,F_t\right).
\end{equation}
其中，$\widehat C_f$ 是频道 $f$ 的定位区域，它包含根据已有检测结果仍不能排除的所有源位置；$\mathcal O_f^+$、$\mathcal O_f^-$ 分别保存该频道实际收到信号和未收到信号的位置。$P_t$ 保存已完成扫描的位置，$F_t$ 保存计划前往但尚未扫描的检测点，以下将这类点简称为待扫描点。记真实源位置为 $p_f$，定位多边形的顶点集为 $V_f$，其最小外接圆沿用第二问的叫法，简称包围圆，记为 $B(z_f,r_f)$。讨论单个源时，简记 $p=p_f$、$x=x_t$、$z=z_f$、$r=r_f$。

\textbf{（3）根据检测结果更新定位区域。}记 $u(\alpha)=(\cos\alpha,\sin\alpha)$，三角函数中的角度统一换算为弧度。设备返回的示向度保留两位小数，因此计算定位区域时采用 $\bar\varepsilon=1.005^\circ$：在题设 $1^\circ$ 的误差范围外多留 $0.005^\circ$，以容纳四舍五入造成的差异。首次检测位置为 $x$、示向度为 $\theta$ 时，先取三角形
\[
\mathcal T(x,\theta)=\operatorname{conv}\left\{x,\quad
x+1500u(\theta)\pm1500\tan\bar\varepsilon\,u(\theta+\pi/2)\right\},
\]
再与场地圆的外切正64边形求交集。三角形包含半径1500 m、方向在示向度两侧 $\bar\varepsilon$ 内的扇形，外切多边形包含整个场地圆，因此近似只会扩大可能范围，不会排除真实源。后续每测得一次示向度，就按第一问的方法增加两个半平面约束，缩小定位区域。这里采用三角形与64边形计算；第二问采用128边形，具体计算不同，但都是用包含原区域的多边形进行近似。

在位置 $x^-$ 检测而未收到信号，说明若该频道有源，它与该点的距离至少为1000 m。先从当前定位区域中排除与该点相距小于1000 m的部分，再取包含所得区域的最小凸多边形，即凸包，以便继续用凸多边形计算；计算中把排除半径内缩到999.999 m，以避免舍入误差影响边界判断。同一频道收到信号的位置 $x^+$ 与未收到信号的位置 $x^-$ 还满足
\begin{equation}\label{eq:q3-3}
\|p-x^+\|\le R_f<\|p-x^-\|
\ \Longrightarrow\ 
2(x^--x^+)^Tp<\|x^-\|^2-\|x^+\|^2.
\end{equation}
将严格不等式放宽为包含边界的半平面，即将“$<$”改为“$\le$”，再将同一频道收到信号和未收到信号的记录逐一配对，用所得约束缩小定位区域。尚未发现该源时，先保存未收到信号的记录，发现后再使用；一次未收到信号并不代表该频道没有源。每次更新后重新计算包围圆。如果定位区域为空，说明检测结果或数值计算之间出现矛盾，不能自行扩大区域后继续。收到近场反馈时，则在当前位置进行光学精确定位和清除。

\Needspace{30mm}
\subsubsection{模型的建立}\label{sec:q3-2}

\textbf{（1）决策内容。}根据当前已知信息 $S_t$，选择下一项任务、检测频道和执行位置。任务分为两类：一类是扫描未知频道，寻找尚未发现的源；另一类是对已发现的源继续测量示向度或执行清除。以 $\mu$ 表示选择动作的方法；每执行一项任务，就根据实际结果更新 $S_t$，再决定下一步。

\textbf{（2）目标函数。}目标是在完成全部清除并确认没有遗漏的条件下，减少任务总耗时。记实际移动总距离为 $L_{\rm move}$，切换频道次数为 $N_{\rm switch}$，无线电检测次数为 $N_{\rm measure}$，光学精确定位尝试次数为 $N_{\rm clear}$，成功清除次数为 $K$。移动速度为5 m/s，切换频道、无线电检测、光学精确定位和清除分别耗时1、5、3和2 s，因此
\begin{equation}\label{eq:q3-2}
T=\frac{L_{\rm move}}{5}+N_{\rm switch}+5N_{\rm measure}+3N_{\rm clear}+2K.
\end{equation}
$N_{\rm clear}$ 包含未找到源的光学精确定位尝试，$K$ 只计实际清除成功的次数，清除后设备仍停留在原检测频道。下文选择路线和检测点时使用的是耗时估计，实际成绩仍按式\eqref{eq:q3-2}统计；程序在计算机上运行所需的时间不计入其中。

\textbf{（3）约束条件。}题目给出的接收、示向度误差和清除距离条件，转化为以下可根据检测结果检查的要求。

\textbf{信息与定位约束}。选择动作只能依据 $S_t$ 中当前已知的信息；更新定位区域时，须保留所有仍可能是源的位置，即保证 $p_f\in\widehat C_f$。为比较候选动作而假设的源位置和示向度，只用于估计耗时，不能用于修改实际定位区域。

\textbf{下一检测点的可靠接收约束}。为缩小已发现源的定位区域而专门选择下一检测点时，沿用第二问的顶点检验，令
\begin{equation}\label{eq:q3-reception}
Q_f=\left\{q:\ \|q-a\|\le1000,\quad a\in V_f\right\}.
\end{equation}
由于多边形内任一点均为顶点的凸组合，$q\in Q_f$ 保证该点到真实源的距离不超过1000 m，即接收半径的下限。选点时还要排除该频道已经检测过的位置，因为同一位置重复检测不能消除固定误差。这是保证下一次定位检测能够收到信号的充分条件。扫描未知频道，以及在途中顺便检测其他已发现源时，并不要求预先满足这一条件，因此仍可能收到无信号的结果。

\textbf{清除距离约束}。收到近场反馈时，可在当前位置进行光学精确定位和清除；否则，只有清除位置到定位区域内每个可能源位置的距离都不超过20 m，才能保证清除成功。需要专门选择清除位置时，先要求包围圆半径 $r_f\le19.5$ m，再取
\begin{equation}\label{eq:q3-11}
\Gamma_f=\bigcap_{a\in V_f}B(a,R_c),\qquad R_c=20-10^{-5}\ \text{m}.
\end{equation}
$\Gamma_f$ 是本文用于选择清除位置的范围，由多个圆盘的交集组成，因此仍是凸集。这里将20 m减去 $10^{-5}$ m，是为了避免数值舍入影响边界判断。任取 $q\in\Gamma_f$ 和 $p=\sum_a\lambda_a a\in\widehat C_f$，由 $\lambda_a\ge0$、$\sum_a\lambda_a=1$，有
\[
\|q-p\|\le\sum_a\lambda_a\|q-a\|\le R_c<20.
\]
因此，只要真实源仍在定位区域内，从该位置到源的距离就小于20 m。确定清除位置前，仍逐个检查它到多边形顶点的距离；只有设备实际反馈清除成功后，才将该源记为已清除。

\textbf{场地覆盖与任务完成约束}。只要还有未知频道，已扫描位置和待扫描点的接收范围就应共同覆盖整个场地，即
\begin{equation}\label{eq:q3-cover-invariant}
\Omega\subseteq\bigcup_{x\in P_t\cup F_t}B(x,1000).
\end{equation}
这一条件保证按照当前计划扫描后，能够发现场地内尚未发现的源。但一个位置即使已经列入计划，只要没有实际扫描，就不能据此判断是否有源。只有已完成扫描的位置 $P_t$ 的接收范围已经覆盖整个场地，才能将仍未发现源的频道确认为没有源；另一种情况是已经发现16个源，此时可根据题目给出的数量上限确认其余频道没有源。已发现但尚未清除的源仍需继续处理。记 $\tau$ 为任务结束时刻，要求
\begin{equation}\label{eq:q3-finish}
\mathcal U_\tau=\varnothing,\qquad
\mathcal D_\tau\setminus\mathcal K_\tau=\varnothing,\qquad
|\mathcal K_\tau\cup\mathcal A_\tau|=20.
\end{equation}
任务结束时，还要核对程序记录的清除数与设备确认的清除成功次数是否一致，再正常结束任务。发现全部源和清除全部源是两个不同条件，不能互相替代。

\textbf{（4）模型汇总。}在检测误差符合题目要求、定位区域包含真实源的条件下，将上述要求归纳为
\begin{equation}\label{eq:q3-model}
\begin{aligned}
\min_{\mu}\quad &T(\mu)\\
\text{s.t.}\quad &\mu\text{ 仅依据当前已知信息 }S_t\text{ 选择动作},\\
&p_f\in\widehat C_f,\qquad f\in\mathcal D_t\setminus\mathcal K_t,\\
&q\in Q_f,\qquad q\text{ 为专门安排定位检测所选的位置},\\
&\text{清除位置具有近场反馈，或对全部顶点满足20 m距离条件},\\
&\text{未知频道的扫描计划满足式}\eqref{eq:q3-cover-invariant},\\
&\text{任务结束时的状态满足式}\eqref{eq:q3-finish}.
\end{aligned}
\end{equation}
式\eqref{eq:q3-model}给出了要减少的耗时和必须满足的条件。由于后续检测结果尚未知，不能在出发前一次算出全部动作。接收半径下限和20 m清除距离来自题目；顶点检验和场地覆盖条件用于保证动作有效。19.5 m的半径条件以及其他略微缩小距离、扩大误差范围的处理，只是为避免数值计算误差，不改变题设参数。

\textbf{（5）目标的分解。}真实源的位置未知，后续检测结果也无法提前确定，因此不能直接算出每一种完整方案的实际耗时 $T$。本文把下一步的选择分为三件事：先做哪项任务、到哪里检测、到哪里清除。分别生成有限个候选方案，再用以下计算方法进行比较。

\emph{任务顺序。}将剩余待扫描点和已发现但未清除的源合并为任务集合 $\mathcal T_t$。安排路线时，暂用各源的包围圆心代表其位置，所有任务的代表位置记为 $y_1,\ldots,y_n$，令 $y_0=x_t$。若按次序 $\pi$ 访问这些位置，且不要求返回原点，则路线长度为
\begin{equation}\label{eq:q3-route-length}
L_{\rm route}(\pi)=\|y_{\pi_1}-y_0\|+
\sum_{j=1}^{n-1}\|y_{\pi_{j+1}}-y_{\pi_j}\|.
\end{equation}
路线不要求返回原点，选择候选路线中 $L_{\rm route}$ 最小者。该长度只对应当前代表点，与式\eqref{eq:q3-2}中的实际累计距离 $L_{\rm move}$ 不同；实际检测位置、清除位置和后来新增的任务仍会随检测结果改变。

\emph{检测点选择。}对待定位源，在当前定位区域内选取若干假想源位置 $g_j$，赋予权重 $w_j$；再选取若干可能的示向度误差 $\eta_k$，赋予权重 $\omega_k$。两组权重均非负，且各自之和为1。记 $H_{jk}(q)$ 为假设源位于 $g_j$、在候选点 $q$ 检测时出现误差 $\eta_k$ 后，继续定位和清除所需的估计时间。$\mathcal Q_{\rm feasible}$ 为去除重复位置并通过接收距离检验的候选检测点集合。定义
\begin{equation}\label{eq:q3-probe-score}
J(q)=\frac{\|q-x\|}{5}+5+
\sum_j w_j\sum_k\omega_kH_{jk}(q),
\qquad q^*\in\arg\min_{q\in\mathcal Q_{\rm feasible}}J(q).
\end{equation}
三项分别是到达检测点的时间、本次检测时间和后续所需时间的估计。假想源位置和 $H_{jk}$ 的计算见模型求解中的式\eqref{eq:q3-probe-future}。位置与误差样本的权重用于构造候选点评分，$J$ 表示后续费用的估计。

\emph{清除位置选择。}若下一项任务的位置为 $y$，就在能够保证清除成功的范围 $\Gamma_f$ 内生成有限个候选位置，以
\begin{equation}\label{eq:q3-through-clear}
G(q)=\|x-q\|+\|q-y\|
\end{equation}
最小者作为清除位置；没有下一项任务时，只比较 $\|x-q\|$。这样既考虑走到清除位置的路程，也考虑清除后去往下一项任务的路程，减少折返。这三个计算分别用于选择当前较省时的动作，不能据此认为整局任务已经达到所有可能方案中的最短时间。

\Needspace{30mm}
\subsubsection{模型的求解}\label{sec:q3-3}

求解分为三部分：布置检测点，保证能够发现尚未发现的源；安排任务先后顺序，减少来回移动；对每个已发现的源继续测量示向度，缩小定位区域，并选择清除位置。每次只执行下一项任务，再根据实际检测结果调整计划。以下结合示意图说明具体做法。

\textbf{（1）初始扫描}先在原点扫描20个频道，再在半径 $a=999$ m的圆周上等间隔布置七个检测点：
\begin{equation}\label{eq:q3-4}
F_0=\left\{a\bigl(\cos(2k\pi/7),\sin(2k\pi/7)\bigr):k=0,\ldots,6\right\}.
\end{equation}
初始布局包含原点与圆周上的七个检测点，共八个位置，但不要求每局都走完八点。原点扫描用于提前发现附近的源；圆周上的七个检测点，其接收范围已经能够覆盖整个场地。

对任意 $p\in\Omega$，令 $r=\|p\|\in[0,1800]$。至少有一个环上点 $x_k$ 与 $p$ 的极角相差不超过 $\pi/7$，于是
\begin{equation}\label{eq:q3-5}
\begin{aligned}
\min_{x_k\in F_0}\|p-x_k\|^2
&\le r^2+a^2-2ar\cos(\pi/7)\\
&\le \max\{a^2,1800^2+a^2-3600a\cos(\pi/7)\}\\
&=998001<1000^2.
\end{aligned}
\end{equation}
第二步利用右侧关于 $r$ 的凸性，最大值在区间端点取得。因此，每个可能的源都至少能被一个环上检测点接收。到达待扫描点后，只扫描仍未知的频道，已发现源按其定位需要另行检测；已经发现16个源后，就停止检测剩余未知频道。图\ref{fig:q3-coverage}画出了这一关系：七个接收圆盘共同盖住场地，任意源都能找到一个极角足够接近的环上检测点。原点扫描先获取起点附近的信息，后续扫描再判断其余未知频道是否有源。

\begin{figure}[!htbp]
\centering
\QThreeCoverageFigure
\caption{原点扫描与七个环上检测点的覆盖布局。左图按实际距离比例绘制；右图说明任意可能源与最近环上检测点的角度关系，用于解释式\eqref{eq:q3-5}。}
\label{fig:q3-coverage}
\end{figure}


\textbf{（2）待扫描点的删除与调整。}依据式\eqref{eq:q3-cover-invariant}，先尝试删除距当前位置最远的待扫描点，再逐个检查较近的点。只有其余待扫描点与已扫描位置的接收范围仍能覆盖整个场地，才接受删除。每删一个点，都要按新的布局重新检查，不能分别判断多个点可以删除后再将它们同时删除。若 $P_t$ 的接收范围已经覆盖整个场地，则将仍未知的频道确认为没有源。

检查是否覆盖整个场地时，先检查场地圆周是否全部被接收圆盘覆盖，再检查各接收圆在场地内部的边界圆弧是否被其他接收圆盘覆盖，后一步用于排除内部仍有未覆盖区域的情况。具体通过计算被覆盖的圆弧区间完成。计算时将场地半径扩大到1800.001 m，将保证接收的半径缩小到999.99 m，以容纳数值误差；不能只检查有限个网格点就认为整个场地已被覆盖。

调整某个待扫描点的位置时，先找出当前路线中它前后的两个任务位置。候选位置包括：前后两点所连线段上距原检测点最近的位置（即线段上的投影点）、原检测点与该投影点的中点，以及附近至多三个已发现源的包围圆心和原检测点到这些圆心的中点；没有后一个任务时，还尝试直接使用前一个任务的位置。只有新位置能使经过该点的路程缩短超过1 m，且调整后的接收范围仍覆盖整个场地，才接受调整。依次检查剩余待扫描点，每次都根据更新后的布局重新计算。只有到达并实际完成扫描的位置，才能加入 $P_t$。

\textbf{（3）任务访问次序的选择。}以式\eqref{eq:q3-route-length}为目标，共尝试40条初始路线，并逐条调整访问次序。前 $n$ 条分别从不同任务开始，之后每一步都选择距离最近的未访问任务，直至补齐路线；其余初始路线由目前最短的路线交换2至4次任务位置得到，交换规则固定，相同输入可以重复得到相同结果。对每条初始路线，先尝试将一段路线的访问次序反转，即2-opt调整；如果不能缩短路程，再尝试把单个任务移到其他位置，包括移到路线末尾。按上述先后顺序反复尝试，每条路线最多调整50轮，最后保留所有尝试中的最短路线。任务最多有七个待扫描点加16个已发现源，因此40条初始路线足以分别尝试每个任务作为第一个任务。

选定路线后，只执行第一项任务，并保留第二项任务的位置供清除落点选择使用。一次检测可能改变源的位置范围或发现新源，因此随后重新排序。图\ref{fig:overall-framework}下方的 2-opt 示意展示局部调整：替换两条连接边并反转中间片段，保留使路线缩短的结果。




\textbf{（4）候选检测点的生成。}对待处理源，若 $\|z-x\|<10^{-7}$ m，认为机器狗已经位于包围圆心，直接在此测量示向度；否则令 $e=(z-x)/\|z-x\|$ 表示从当前位置指向包围圆心的单位向量，$e_\perp=(-e_y,e_x)$ 表示与它垂直的单位向量。候选检测点取为
\begin{equation}\label{eq:q3-probe-candidates}
\mathcal Q=\{x\}\cup
\left\{z+r(\alpha e+\beta e_\perp):
\begin{array}{l}
\alpha\in\{-0.7,-0.35,0,0.35,0.7\},\\
\beta\in\{-0.4,-0.15,0,0.15,0.4\}
\end{array}\right\}.
\end{equation}
筛选前至多有26个候选位置。只保留到定位多边形每个顶点的距离都小于999.9 m，且该频道未在此处检测过的位置，得到 $\mathcal Q_{\rm feasible}$。999.9 m比接收半径下限1000 m略小，以避免数值误差影响接收保证。本问沿用第二问先生成候选点、再检验接收距离的思路，具体位置则根据包围圆半径确定。如图\ref{fig:q3-probe}，候选点沿 $e$ 和 $e_\perp$ 两个方向分布：前后位置影响接近源的距离，横向位置影响两次示向方向的交会角。先排除不能保证收到信号的位置，再用式\eqref{eq:q3-probe-score}比较走到各点的时间，以及检测后预计还需花费的时间。

\begin{figure}[!htbp]
\centering
\QThreeProbeFigure
\caption{候选检测点的生成与距离筛选示意。}
\label{fig:q3-probe}
\end{figure}


\textbf{（5）候选检测点的耗时估计。}从定位多边形的一个顶点向其他不相邻顶点连线，将其分成若干三角形。对每个三角形 $(a,b,c)$，在
\[
\frac{a+b+c}{6}+\frac a2,\qquad
\frac{a+b+c}{6}+\frac b2,\qquad
\frac{a+b+c}{6}+\frac c2
\]
三个位置取样，每个点赋予该三角形面积占总面积的三分之一权重，得到 $g_j,w_j$；如果区域退化为线段或点，则直接取顶点并赋予相同权重。三个假设的示向度误差及其权重取为
\[
\eta_k\in\{-\sqrt{3/5}\,\varepsilon,0,\sqrt{3/5}\,\varepsilon\},
\qquad (\omega_1,\omega_2,\omega_3)=(5/18,4/9,5/18),\quad \varepsilon=1^\circ.
\]
假设源在 $g_j$，在候选点 $q$ 测得的示向度为 $\arg(g_j-q)+\eta_k$。先复制当前定位区域，再加上该示向度对应的楔形约束，计算假设收到这一结果后区域会如何变化。这里的楔形半角取 $1^\circ$；实际检测后的区域更新仍采用 $1.005^\circ$。在假设得到的新区域中，以最远两个顶点的中点 $\widetilde z_{jk}$ 为圆心，以到所有顶点的最大距离 $\widetilde r_{jk}$ 为半径，构造一个包含整个区域的圆，用于估计后续耗时。这样不会将直径的一半误当成能覆盖整个区域的圆半径。

估计耗时时，允许安排清除的包围圆半径上限取 $r_c=19.5$ m。若 $\|q-g_j\|\le5$ m，则预计会收到近场反馈，可以直接进行光学精确定位和清除，两项共需5 s，取 $H_{jk}=5$ s；否则令
\[
d_{jk}=\|q-\widetilde z_{jk}\|,\qquad
b_{jk}=\|g_j-\widetilde z_{jk}\|,\qquad
\ell_{jk}=\max\{0,\log_2(\max\{\widetilde r_{jk},r_c\}/r_c)\},
\]
按下式计算假设得到该检测结果后的耗时：
\begin{equation}\label{eq:q3-probe-future}
H_{jk}(q)=
\begin{cases}
\displaystyle\frac{[d_{jk}-(20-\widetilde r_{jk})]_+}{5}+5,
&\widetilde r_{jk}\le r_c,\\[5pt]
\displaystyle\frac{d_{jk}+[b_{jk}-r_c]_+}{5}
+6\ell_{jk}+0.1\widetilde r_{jk}+5,
&\widetilde r_{jk}>r_c,
\end{cases}
\end{equation}
其中 $[a]_+=\max\{a,0\}$。第一种情形估计走到能够保证清除成功的位置，再完成光学精确定位和清除所需的时间；第二种情形还考虑继续接近源、通过检测和切换频道缩小定位区域所需的时间，并按上述近似圆的半径 $\widetilde r_{jk}$ 增加一项估计耗时。代入式\eqref{eq:q3-probe-score}计算各候选点的 $J(q)$，取最小者。这些假设结果只用于比较候选点，实际定位区域必须在收到真实检测结果后才能更新。

\textbf{（6）在包围圆心检测，继续缩小定位区域。}对同一源专门安排的定位检测达到4次，或当前包围圆半径小于22 m时，改到包围圆心测量示向度；没有新的可行候选点时，也采用包围圆心。若 $\widehat C_f\subseteq B(z,r)$，在 $z$ 测得符合误差范围的示向度后，新的定位区域包含于半角 $\bar\varepsilon$、半径 $r$ 的扇形，因此其包围圆半径满足
\begin{equation}\label{eq:q3-center-contraction}
r_{\rm new}\le\frac{r}{2\cos\bar\varepsilon}.
\end{equation}
经过扇形顶点和圆弧两个端点的圆，可以覆盖整个扇形，由此得到这一半径上界。每次在包围圆心检测后，都检查新的包围圆半径是否满足该关系，比较时允许 $10^{-3}$ m的数值差异。如果在圆心收不到信号，或新的半径不满足该关系，则报错并停止。如此继续检测，逐步缩小定位区域，直到满足清除距离条件。

图\ref{fig:q3-contraction}给出一个计算例子：若开始在包围圆心检测时，包围圆半径为160 m，按上述关系依次计算，半径上界约为80.01、40.01、20.01和10.01 m。只要每次检测误差符合要求，且更新后的定位区域始终包含真实源，至多四次即可使半径上界低于19.5 m。若提前满足清除条件或收到近场反馈，就可以提前清除。这里展示的是公式给出的半径变化，不是实际检测数据。

\begin{figure}[!htbp]
\centering
\QThreeContractionFigure
\caption{在包围圆心检测后定位区域的变化及半径上界。}
\label{fig:q3-contraction}
\end{figure}


\textbf{（7）清除位置的选择。}每轮先检查当前位置到某个源的定位区域所有顶点是否均不超过20 m，若满足，就立即进行光学精确定位和清除。需要另外选择位置时，先检查包围圆半径是否不超过19.5 m，再按式\eqref{eq:q3-11}计算允许清除的位置集合 $\Gamma_f$。如果没有下一项任务，就选择 $\Gamma_f$ 中距当前位置最近的点；如果下一项任务的位置为 $y$，分别求 $\Gamma_f$ 中距 $x$、$y$ 最近的点 $q_x,q_y$，在连线 $q_xq_y$ 上等间隔取21点作为候选。再检查线段 $xy$ 是否穿过 $\Gamma_f$，若存在公共线段，则把该线段的中点也加入候选。比较式\eqref{eq:q3-through-clear}，取总路程最短者。这样不必每次都走到包围圆心，但选择的位置仍须满足到全部顶点的距离条件。

\textbf{（8）利用当前位置补充检测。}完成一次定位检测或清除后，先试算：如果在当前位置扫描，是否可以因此删除至少一个待扫描点。若可以，就在此处实际扫描仍未知的频道，然后更新已扫描位置和后续计划，用这次扫描替代原定检测点。

若不在当前位置扫描未知频道，再考虑是否值得检测其他已发现源所在的频道。对已经满足清除所需半径、已在当前位置检测过，或满足 $\|x-z_f\|-r_f>1500$ 的源，不再测量示向度。对其余频道，先假设源位于其包围圆心，计算相应示向度，再求加入该角度约束后的定位区域。如果新区域的包围圆半径小于原半径的80\%，或已满足清除所需半径，就进一步计算：以新区域的外接矩形中心为圆心，以到区域各顶点的最大距离为半径，记为 $r^{\prime}_{\rm box}$。当
\begin{equation}\label{eq:q3-opportunity}
r'_{\rm box}\le19.5\quad\text{或}\quad r_f-r'_{\rm box}>30\ \text{m}
\end{equation}
时，再实际检测一次。到达待扫描点并完成未知频道扫描后，也按上述包围圆半径缩小至原来80\%以下或满足清除所需半径的条件，决定是否补充检测已发现源的频道。这里的近似计算都只用于决定要不要检测，实际定位区域仍由真实检测结果更新。

上述规则只决定是否值得再检测一次。没有实际检测时，不能把该位置记为“未收到信号”，也不能因此排除可能的源位置；没有扫描过未知频道的位置，也不能计入已完成的场地覆盖。每到一个检测点，优先检测设备当前所在的频道，以减少切换频道的次数。

\textbf{（9）求解步骤与结束检查。}每发现或清除一个源，就检查已发现源数和实际完成的场地覆盖情况。满足式\eqref{eq:q3-finish}后，核对清除成功次数并结束任务。程序最多重复选择下一项任务600轮，防止出现异常后一直运行；场地覆盖和各源定位区域也分别检查。如果超过600轮、检测结果互相矛盾、清除失败，或设备没有返回明确结果，就记录任务中断或失败，不能算作完成全部清除。图\ref{fig:q3-flow}给出执行流程，算法\ref{alg:q3-search}列出具体步骤。

\begin{figure}[!htbp]
\centering
\QThreeFlowFigure
\caption{多源搜索与清除模型的求解流程。}
\label{fig:q3-flow}
\end{figure}

\FloatBarrier


\Needspace{70mm}
\par\noindent\begin{minipage}{\linewidth}
\refstepcounter{algorithm}\label{alg:q3-search}
\textbf{算法\thealgorithm：多源搜索与清除模型的求解}
\small
\begin{enumerate}[label=\arabic*.,ref=\arabic*,leftmargin=2em,itemsep=2pt,topsep=4pt]
  \item 初始化20个频道与圆周上的七个检测点，在原点扫描，并根据实际检测结果建立定位区域。
  \item\label{step:q3-finish-check} 检查已发现源数和实际扫描范围，判断哪些频道可以确认没有源；若满足式\eqref{eq:q3-finish}，核对清除成功次数并结束任务。
  \item 清除在当前位置满足距离条件的源；检查哪些待扫描点可以删除，或可以调整位置以缩短路程。
  \item 将待扫描点与已发现源一起安排，通过40条初始路线及其后续调整选择下一项任务，并保存再下一项任务的位置。
  \item 若为扫描任务，到达检测点后扫描未知频道，并按预计能否缩小区域来决定是否补测已发现源；否则，清除指定的源，或到选出的检测点测量示向度，必要时改到包围圆心检测。
  \item 判断是否在当前位置扫描未知频道，或补充检测其他已发现源；根据实际结果更新定位区域、清除记录和已扫描位置。
  \item 如果尚未完成任务，且未超过程序限制、未出现设备返回异常，则回到步骤\ref{step:q3-finish-check}；否则保存完成或中断的原因。
\end{enumerate}
\end{minipage}\par


场地覆盖的证明说明按计划实际扫描能够发现未知源，在包围圆心继续检测则提供了逐步缩小定位区域的方法。实际完成率和耗时由下节实验评价。
%
\Needspace{60mm}
\subsubsection{结果分析}\label{sec:q3-4}


\textbf{（1）实验设置与评价指标。}本地实验采用本地模拟器生成的400个场景，每场含10至16个干扰源，位置在半径1800 m的圆形场地内按面积均匀分布。

设第 $i$ 场共有 $N_i$ 个干扰源，完成任务的总耗时为 $T_i$。评价时先将每场耗时除以该场干扰源数，再将各场结果取算术平均，同时报告平均每局耗时：
\begin{equation}\label{eq:q3-17}
\overline q=\frac1n\sum_{i=1}^{n}\frac{T_i}{N_i},
\qquad
\overline T=\frac1n\sum_{i=1}^{n}T_i.
\end{equation}
其中，$\overline q$ 的单位为秒/源，$\overline T$ 的单位为秒/局。总耗时包括清除最后一个干扰源后，为确认没有遗漏而进行的检测与移动，也包括光学精确定位失败时的耗时；此外记录完成全部清除的局数和实际清除的干扰源数。$\overline q$ 先对每场单独计算，再取平均，与直接将全部时间除以全部干扰源数的结果不同。

%
\par\medskip\Needspace{4\baselineskip}
\textbf{（2）任务完成情况。}本问方案在400个本地场景中全部完成清除，共清除5209个干扰源，每个源的平均耗时为243.991 s/源，平均每局耗时为3119.627 s/局。直接调用模拟器进行演练的20局也全部完成，共清除277个干扰源，每个源的平均耗时为225.623 s/源，平均每局耗时为3074.968 s/局。两批记录中的光学精确定位尝试均成功，结果见表\ref{tab:q3-current-results}。

\Needspace{40mm}
\begin{center}
\begin{minipage}{\linewidth}
\centering\small
\captionof{table}{原点扫描方案的任务完成情况与耗时}
\label{tab:q3-current-results}
\vspace{4pt}
\begin{tabularx}{\linewidth}{@{}Xrrrrr@{}}
\toprule
数据来源 & \shortstack{全部清除\\局数} & 清除源数 & \shortstack{平均耗时\\(s/源)} & \shortstack{平均耗时\\(s/局)} & \shortstack{光学精确定位\\失败次数} \\
\midrule
本地场景 & 400/400 & 5209 & 243.991 & 3119.627 & 0 \\
模拟器演练 & 20/20 & 277 & 225.623 & 3074.968 & 0 \\
\bottomrule
\end{tabularx}
\end{minipage}
\end{center}

这些完成记录说明，在所测试的场景中，机器狗能够通过扫描发现干扰源，用实际测得的示向度缩小定位区域，在满足距离条件的位置完成清除，并在结束前确认没有遗漏。
%
\input{sections/q3/04_training_and_practice}%
\input{sections/q3/05_trajectories_and_discussion}%

本问的圆盘覆盖和包围圆心测向都利用了全向接收条件。下一问保留滚动规划的任务组织方式，进一步处理发射方向未知时的覆盖与定位问题。
